\documentclass[aspectratio=169,usenames,dvipsnames]{beamer}
\beamertemplatenavigationsymbolsempty
% \documentclass{beamer}

% UNCOMMENT FOR HANDOUTS
% \usepackage{pgfpages}
% \pgfpagesuselayout{2 on 1}[landscape,letterpaper,border shrink=5mm]
% \setbeameroption{show notes}

\mode<presentation>
\usetheme{Boadilla}
\usecolortheme{spruce}
\usepackage{../../LizStyle,ulem}
\usepackage{color}
\usepackage[color,matrix,arrow]{xy}
% \input xy
% \xyoption{all}
\usepackage{tikz}
\usepackage{tikz-cd}
\tikzset{commutative diagrams/diagrams={ampersand replacement=\&}}

\AtBeginSection{\frame{\sectionpage}}


% ==========Command for putting comments and removing them=======
% Visible:
\newcommand{\InstructorNotes}[1]{\textit{\color{orange}#1}}
\newcommand{\InstructorList}[1]{
\begin{itemize}
	\color{orange}
	#1
\end{itemize}
}
% Removed:
% \renewcommand{\InstructorNotes}[1]{}
% \renewcommand{\InstructorList}[1]{}


\newcommand{\bvfill}{\vskip0pt plus 1filll}

\usepackage{multimedia}
\usepackage{graphbox}

% \usepackage{algpseudocode}


\newcommand{\Reeb}{\textbf{Reeb}}
\newcommand{\Rtopc}{\textbf{Rtopc}}
% \newcommand{\RR}{\mathcal{R}}

\newcommand{\Set}{\textbf{Set}}
\newcommand{\Vect}{\textbf{Vect}}
\newcommand{\Top}{\textbf{Top}}
\newcommand{\Int}{\textbf{Int}}


\graphicspath{ {Figures/} {../Figures/}{../../Figures/} }



\begin{document}
%\pdfoptionpdfminorversion 5
%\logo{\includegraphics[height=1.5cm]{dukeshield.pdf}}
\title[]{Ch 9.1-2: Maximum Margin Classifier and Support Vector Classifier}
\subtitle{Lecture 23 - CMSE 381}
\date{Mon, Apr 11, 2022}



\institute[MSU-CMSE]{Michigan State University\\
			::\\
          Dept of Computational Mathematics, Science \& Engineering}
\author[Dr.~Munch]{Prof.~Elizabeth Munch}

\begin{frame}
 \titlepage


\end{frame}
%--------------------------
\begin{frame}[c]{Announcements }

\begin{columns}
\column{.45\textwidth}
\centering

	\textbf{Last time:}
	\begin{itemize}
	\item Ch 8: Decision Trees
	\end{itemize}

	\vspace{.2in}

	\textbf{This lecture:}
	\begin{itemize}
	\item Maximal Margin Classifier
	\item Support Vector Classifier
\end{itemize}


\column{.45\textwidth}
\centering

\textbf{Announcements:}
	\begin{itemize}
	\item Mon 4/11: 9.1-2
	\item Weds 4/13: 9.3-4
	\item Fri 4/1: Lab: 9.6
	\begin{itemize}
		\item HW \#9 Due Friday
	\end{itemize}

	\end{itemize}
\end{columns}

\end{frame}
%--- Next Frame ---%



%=================
\section{Maximal Margin Classifier}
%=================

\begin{frame}[t]{What is a hyperplane?}

\begin{columns}
\column{.45\textwidth}
\centering
\InstructorList{
  \item In $p$-dimensional space, a hyperplane is a flat affine dubspace of dimension $p-1$
	\item Draw example for $p=2,3$
	\item
}


\column{.45\textwidth}
\centering

\end{columns}

\end{frame}
%--- Next Frame ---%
\begin{frame}[t]{Mathematical definition of a hyperplane}

	\begin{equation*}
		\beta_0 + \beta_1X_1 + \beta_2X_2 + \cdots + \beta_p X_p = 0
	\end{equation*}

\begin{columns}
\column{.45\textwidth}
\centering
\InstructorList{
  \item Being on the hyper plane means having a point $X=(X_1,X_2,\cdots,X_p)^T$ satisfying this
	\item Show that when $p=2$, this is just a line
}

\column{.45\textwidth}
\centering

\end{columns}

\end{frame}
%--- Next Frame ---%

\begin{frame}[t]{Hyperplane for $p=2$}

\begin{columns}
\column{.45\textwidth}
\centering
\includegraphics[width = \textwidth, align = c]{Fig9-1}
$1+2X_1 + 3X_2 = 0$
\column{.45\textwidth}
\centering

\end{columns}

\end{frame}
%--- Next Frame ---%
\begin{frame}[t]{There are two sides to every hyperplane}

\begin{columns}
\column{.45\textwidth}
\centering
\begin{equation*}
		\beta_0 + \beta_1X_1 + \beta_2X_2 + \cdots + \beta_p X_p < 0
\end{equation*}
\column{.45\textwidth}
\centering
\begin{equation*}
		\beta_0 + \beta_1X_1 + \beta_2X_2 + \cdots + \beta_p X_p > 0
\end{equation*}
\end{columns}
\bvfill
\centering
\includegraphics[width = .4\textwidth, align = c]{Fig9-1}

\end{frame}
%--- Next Frame ---%


\begin{frame}[t]{Classification Setup}

\begin{columns}
\column{.45\textwidth}
\centering
 Data matrix:
	\begin{equation*}
		X =
		\begin{pmatrix}
			- & x_1^T & -\\
			- & x_2^T & -\\
			 & \vdots & \\
			- & x_n^T & -
		\end{pmatrix}_{n \times p}
	\end{equation*}

	\begin{equation*}
		x_1 =
		\begin{pmatrix}
			x_{11} \\
			\vdots \\
			x_{1p}
		\end{pmatrix}
		,\cdots,
		x_n =
		\begin{pmatrix}
			x_{n1} \\
			\vdots \\
			x_{np}
		\end{pmatrix}
	\end{equation*}

\column{.45\textwidth}
\centering
 Observations in one of two classes, $y_i \in \{-1,1\}$
	\begin{equation*}
		Y =
		\begin{pmatrix}
			y_1 \\
			y_2\\
			\vdots \\
			y_n
		\end{pmatrix}
	\end{equation*}


% \column{.4\textwidth}
% \centering

\end{columns}
\bvfill
\centering
Separate out a test observation
\begin{equation*}
	x^* = (x_1^* \; \cdots \; x^*_p)^T
\end{equation*}
\end{frame}
%--- Next Frame ---%


\begin{frame}[t]{Separating Hyperplane}

\vspace{-.2in}
\begin{align*}
	\beta_0 + \beta_1x_{i1} + \beta_2 x_{i2} + \cdots + \beta_px_{ip} >0 \text{ if } y_i & = 1\\
	\beta_0 + \beta_1x_{i1} + \beta_2 x_{i2} + \cdots + \beta_px_{ip} <0 \text{ if } y_i & = -1\\
\end{align*}
\vspace{-.4in}
\begin{columns}
\column{.4\textwidth}
\centering
\includegraphics[width = \textwidth, align = c]{Fig9-2a}
\column{.45\textwidth}
\centering
\InstructorList{
  \item Strict requirement
	\item No points can be on the wrong side of the line
}

\end{columns}

\end{frame}
%--- Next Frame ---%
\begin{frame}[t]{Another way to say it}

\begin{align*}
	\beta_0 + \beta_1x_{i1} + \beta_2 x_{i2} + \cdots + \beta_px_{ip} >0 \text{ if } y_i & = 1\\
	\beta_0 + \beta_1x_{i1} + \beta_2 x_{i2} + \cdots + \beta_px_{ip} <0 \text{ if } y_i & = -1\\
\end{align*}

\bvfill
\centering
For all $i$:
\begin{equation*}
	y_i(\beta_0 + \beta_1x_{i1} + \beta_2 x_{i2} + \cdots + \beta_px_{ip}) > 0
\end{equation*}
\bvfill

\end{frame}
%--- Next Frame ---%
\begin{frame}[t]{Separating hyperplane becomes a classifier }

\begin{columns}
\column{.45\textwidth}
\centering
If you have a separating hyperplane:
\begin{itemize}
	\item Check
	\begin{equation*}
		f(x^*) = \beta_0 + \beta_1x_{1}^* + \beta_2 x_{2}^* + \cdots + \beta_px_{p}^*
	\end{equation*}
	\item If positive, assign $\hat y = 1$
	\item If negative, assign $\hat y = -1$
\end{itemize}
\column{.45\textwidth}
\centering
\includegraphics[width = \textwidth, align = c]{Fig9-2b}
\end{columns}

\end{frame}
%--- Next Frame ---%
\begin{frame}[t]{How do we pick?}

\begin{columns}
\column{.45\textwidth}
\centering
\includegraphics[width = \textwidth, align = c]{Fig9-2a}
\column{.45\textwidth}
\centering
\InstructorList{
  \item
\InstructorNotes{There's lots of separating hyperplanes here, so how do we find a "best" one?}
\item Compute the perpendicular distance from each training observation to a given separating hyperplane
\item Smallest such distance is the \textit{margin}
}

\end{columns}

\end{frame}
%--- Next Frame ---%
\begin{frame}[t]{Maximal margin classifier}

\begin{columns}
\column{.45\textwidth}
\centering
\includegraphics[width = \textwidth, align = c]{Fig9-3}
\column{.45\textwidth}
\centering
\InstructorList{
\item \textit{Maximal margin hyperplane} is the hyperplane with the largest minimum distance to the training observations
\item Classify using it gives you the \textit{maximal margin classifier}
\item Observations that are closest are called \textit{support vectors}
\item If points moved slightly then MMH would move
\item No effect if any other points moved
}

\end{columns}

\end{frame}
%--- Next Frame ---%

\begin{frame}[t]{Mathematical Formulation}

\begin{columns}
\column{.65\textwidth}
\centering
\includegraphics[width = \textwidth, align = c]{Eq9-9}
\column{.35\textwidth}
\centering
\InstructorList{
  \item Last eq forces points to be on the correct side of the hyperplane if $M>0$
	\item Making $M$ as big as possible is the maximal margin hyperplane
}

\end{columns}

\end{frame}
%--- Next Frame ---%

\begin{frame}[t]{But what if....}

\begin{columns}
\column{.45\textwidth}
\centering
\includegraphics[width = \textwidth, align = c]{Fig9-4}
\column{.45\textwidth}
\centering
\InstructorList{
  \item
\InstructorNotes{...there is no separating hyperplane}
\item
\InstructorNotes{Means that there is no solution for $M>0$}
}


\end{columns}

\end{frame}
%--- Next Frame ---%

\begin{frame}[t]{Sensitivity to new points}
\InstructorNotes{Even if there is a spearating hyperplane might not be good}
\begin{columns}
\column{.4\textwidth}
\centering
\includegraphics[width = \textwidth, align = c]{Fig9-5a}
\column{.4\textwidth}
\centering
\includegraphics[width = \textwidth, align = c]{Fig9-5b}

\end{columns}
\footnotesize
\InstructorList{
  \item Add a blue observation on the right
	\item Dramatic shift in maximal margin hyperplane
}

\end{frame}
%--- Next Frame ---%
%==========
\section{Support Vector Classifier}
%==========
\begin{frame}[t]{Basic idea}

\begin{columns}
\column{.45\textwidth}
\centering
\InstructorList{
  \item Be ok with having a classifier that isn't quite perfect
	\item Aim for greater robustness to individual observations
	\item Better classification of most of the training observations
	\item Result is a support vector classifier
	\item Soft margin classifier
}

\column{.45\textwidth}
\centering

\end{columns}

\end{frame}
%--- Next Frame ---%
\begin{frame}[t]{Soft margin}

\begin{columns}
\column{.45\textwidth}
\centering
\includegraphics[width = \textwidth, align = c]{Fig9-6a}
\InstructorNotes{Some points on wrong side of margin}
\column{.45\textwidth}
\centering
\includegraphics[width = \textwidth, align = c]{Fig9-6b}
\InstructorNotes{Some points on wrong side of hyperplane (Misclassified)}
\end{columns}

\end{frame}
%--- Next Frame ---%
\begin{frame}[t]{Math version of SVC}

\begin{columns}
\column{.65\textwidth}
\centering
\includegraphics[width = \textwidth, align = c]{Eq9-12}
\column{.35\textwidth}
\centering
\InstructorList{
  \item $C$ is nonnegative tuning parameter
	\item $M$ is the width of the margin
	\item $\e_1, \cdots, \e_n$ are slack variables allowing observations to go to the other side
}

\end{columns}
\end{frame}
%--- Next Frame ---%

\begin{frame}[t]{What is $\e$?}

\begin{columns}
\column{.65\textwidth}
\centering

\includegraphics[width = \textwidth, align = c]{Eq9-12}
\footnotesize
\InstructorList{
  \item If $\e_i = 0$, then on correct side of margin
	\item If $\e_i>0$ then on the wrong side of margin  (Violated margin)
	\item If $\e_i >1$ then on the wrong side of hyperplane
}
\column{.25\textwidth}
\centering
\includegraphics[width = \textwidth, align = c]{Fig9-6a}
\includegraphics[width = \textwidth, align = c]{Fig9-6b}

\end{columns}

\end{frame}
%--- Next Frame ---%


\begin{frame}[t]{What is $C$?}

\begin{columns}
\column{.65\textwidth}
\centering

\includegraphics[width = \textwidth, align = c]{Eq9-12}
\footnotesize
\InstructorList{
\item Bounds sum of $\e_i$, so controls number \& severity of violating margin (budget)
\item $C=0$ means no violations allowed
\item $C>0$ means at most $C$ observations can be on wrong side of hyperplane
}
\column{.25\textwidth}
\centering
\includegraphics[width = \textwidth, align = c]{Fig9-6a}
\includegraphics[width = \textwidth, align = c]{Fig9-6b}

\end{columns}

\end{frame}
%--- Next Frame ---%
\begin{frame}[t]{Examples messing with $C$}

\centering

\includegraphics[width = .24\textwidth, align = c]{Fig9-7a}
\includegraphics[width = .24\textwidth, align = c]{Fig9-7b}
\includegraphics[width = .24\textwidth, align = c]{Fig9-7c}
\includegraphics[width = .24\textwidth, align = c]{Fig9-7d}

Decreasing $C$ $\rightarrow$


\end{frame}
%--- Next Frame ---%
\begin{frame}[t]{What affects the hyperplane?}

\begin{columns}
\column{.45\textwidth}
\centering

\includegraphics[width = \textwidth, align = c]{Fig9-7d}
\column{.45\textwidth}
\centering
\InstructorList{
  \item Only observations on the margin or violating the margin affect the hyperplane
	\item These observations are called support vectors
	\item Changing other points positions doesn't affect hyperplane found
	\item
}

\end{columns}

\end{frame}
%--- Next Frame ---%

\begin{frame}[t]{Bias-variance tradeoff}

\begin{columns}
\column{.45\textwidth}
\centering
\InstructorList{
  \item Large $C$ $\Rightarrow$ wide margin $\Rightarrow$ many support vectors
	\item This results in low variance but high bias
}

\column{.45\textwidth}
\centering
\InstructorList{
  \item Small $C$ $\Rightarrow$ small margin $\Rightarrow$ few support vectors
	\item This results in high variance but low bias
}

\end{columns}

\bvfill
\centering
\includegraphics[width = .24\textwidth, align = c]{Fig9-7a}
\includegraphics[width = .24\textwidth, align = c]{Fig9-7b}
\includegraphics[width = .24\textwidth, align = c]{Fig9-7c}
\includegraphics[width = .24\textwidth, align = c]{Fig9-7d}

Decreasing $C$ $\rightarrow$
\end{frame}
%--- Next Frame ---%
%=================

\begin{frame}[t]{TL;DR}

\begin{columns}
\column{.45\textwidth}
\centering
\textbf{Maximal Margin Classifier}
\includegraphics[width = \textwidth, align = c]{Eq9-9}
\includegraphics[width = .5\textwidth, align = c]{Fig9-3}
\column{.45\textwidth}
\centering
\textbf{Support Vector Classifier}
\includegraphics[width = \textwidth, align = c]{Eq9-12}
\includegraphics[width = .5\textwidth, align = c]{Fig9-7d}

\end{columns}

\end{frame}
%--- Next Frame ---%

\begin{frame}[t]{Next time}


\begin{columns}
\column{.45\textwidth}
\centering

\textbf{Next time:}

\begin{itemize}
	\item 9.3/9.4
\end{itemize}

\column{.45\textwidth}
\centering

\textbf{This week:}
\begin{itemize}
\item Mon 4/11: 9.1-2
\item Weds 4/13: 9.3-4
\item Fri 4/1: Lab: 9.6
\begin{itemize}
	\item HW \#9 Due Friday
\end{itemize}

\end{itemize}
	\end{columns}

\end{frame}
%--- Next Frame ---%


\end{document}
